% `tex.web` §580: `\fontdimen` is bounded in *both* directions.
%
%     if n<=font_params[f] then …
%     else begin while font_ptr=f do  {the font is the most recently loaded}
%       begin … incr(font_params[f]); … end;
%       if n>font_params[f] then
%         begin print_err("Font "); print_esc(font_id_text(f));
%         print(" has only "); print_int(font_params[f]);
%         print(" fontdimen parameters");
%         help2("To increase the number of font parameters, you must")
%         ("use \fontdimen immediately after the \font is loaded."); error;
%         end
%       end
%
% Three rules, and this engine had none of them. It kept `\fontdimen` in a side
% table that answered every read, so:
%
%   * assigning past the last parameter of a font that is no longer the most
%     recently loaded silently stored a value TeX drops;
%   * reading past it silently returned that value, where TeX raises the same
%     error a second time and reads nought;
%   * and the font's own parameter *count* never grew, so a font could never be
%     promoted.
%
% The count is the part that reaches beyond diagnostics: §1195 measures a math
% family by it. Twenty-two parameters make a symbol font and thirteen an
% extension font, whatever the `.tfm` came with — trip.tex has a macro for
% exactly that promotion, and without it we reported `Insufficient symbol
% fonts` at a point where TeX had got as far as `Insufficient extension fonts`.
\catcode`\{=1 \catcode`\}=2
\errorcontextlines=5
\font\r=cmr10 \font\e=cmex10
\message{<n=\the\fontdimen7\r>}
\fontdimen8\r=1pt
\message{<after=\the\fontdimen8\r>}
\fontdimen8\e=2pt
\message{<last=\the\fontdimen8\e>}
\end
